\appendix
\specialsection{Приложение А. Валидационная выборка видео}
\label{appendix:validation-set}

В приложении приведён полный список видео, использованных для проверки работы
модулей системы. Для каждого видео указаны идентификатор, источник, ключевые
выходы пайплайна и причина включения в выборку.



\begingroup
\scriptsize
\sloppy
\emergencystretch=2em
\setlength{\tabcolsep}{2.0pt}
\renewcommand{\arraystretch}{1.14}
\newcolumntype{L}[1]{>{\raggedright\arraybackslash}p{#1}}

\begin{landscape}

\begin{longtable}{
L{0.035\linewidth}
L{0.235\linewidth}|
L{0.075\linewidth}|
L{0.085\linewidth}
L{0.160\linewidth}
L{0.070\linewidth}
L{0.085\linewidth}|
L{0.035\linewidth}|
L{0.190\linewidth}
}

\caption{Валидационная выборка видео и ключевые результаты работы пайплайна.
Сокращения: A/V/T — audio/video/technical score; A/I — уровни академичности и инструментальности.
В столбцах флагов указаны не только численные триггеры, но и краткая
интерпретация причины срабатывания.
 Для сенсорного флага приведены VLM-факторы: F1--F4 относятся к перегруженности слайда/кадра, U1--U4 — к визуальной неаккуратности и посторонним элементам;
расшифровка факторов приведена в \cref{subsubsec:sensor}, как и пояснения для других флагов в \cref{sec:method-flags}.}
\label{tab:validation-set-full}\\

\toprule
\textbf{ID} &
\textbf{Видео / источник} &
\textbf{A/V/T} &
\textbf{Логик} &
\textbf{Сенсорик} &
\textbf{Интуит} &
\textbf{Эмоция} &
\textbf{A/I} &
\textbf{Зачем включено} \\
\midrule
\endfirsthead

\toprule
\textbf{ID} &
\textbf{Видео / источник} &
\textbf{A/V/T} &
\textbf{Логик} &
\textbf{Сенсорик} &
\textbf{Интуит} &
\textbf{Эмоция} &
\textbf{A/I} &
\textbf{Зачем включено} \\
\midrule
\endhead

\hypertarget{vid:v01}{V01} &
\textit{Предел функции. Определение предела функции ``по Коши'' и ``по Гейне''}. Борис Трушин. {\tiny\url{https://www.youtube.com/watch?v=UzfAt6DoN3E}} &
8.1 / 6.9 / 7.5 &
--- &
--- &
--- &
--- &
4 / 2 &
Базовый «хороший средний» пример: качественная запись, планшетная доска, умеренная академичность, без ожидаемых флагов. \\

\midrule

\hypertarget{vid:v02}{V02} &
\textit{Свойства пределов последовательностей, связанные с неравенствами}. Борис Трушин. {\tiny\url{https://www.youtube.com/watch?v=WTjfi-eqL7E}} &
7.3 / 3.7 / 4.9 &
--- &
$\checkmark$ Монтаж/качество: jerky/min=0.56, video=3.7<5. VLM не основной; единичный U3-пересвет. &
--- &
--- &
4 / 2 &
Похожий на V01 по содержанию и подаче, но с более старым техническим качеством и заметными склейками; проверка видео- и sensor-компонента. \\

\midrule

\hypertarget{vid:v03}{V03} &
\textit{Введение в математический анализ 4. Предел последовательности}. Лекторий ФПМИ. {\tiny\url{https://youtu.be/x8KWbcDxnu0?si=OS7ZcAYwkWozgc1C}} &
7.8 / 7.3 / 7.5 &
--- &
--- &
--- &
--- &
4 / 2 &
Классическая длинная университетская лекция у доски: проверка академического профиля и укрупнённой сегментации при цели «составить общее представление». \\

\midrule

\hypertarget{vid:v04}{V04} &
\textit{Садовничая И.В. Лекция 1 по математическому анализу}. Ёжик в матане. {\tiny\url{https://www.youtube.com/watch?v=JWPstFGjI-g}} &
3.6 / 6.9 / 4.7 &
$\checkmark$ Низкое качество звука: audio=3.6<5. &
$\checkmark$ VLM: U1,U3,U4; 3/5 кадров. Неаккуратный фон, предметы, разводы и следы на доске. &
--- &
--- &
3 / 1 &
Похожий сеттинг на МФТИ, но другая подача и хуже звук; проверка штрафа за аудио и устойчивости профиля подачи к стилю лектора. \\

\midrule

\hypertarget{vid:v05}{V05} &
\textit{Предельно понятное видео о пределе}. Профиматика. {\tiny\url{https://www.youtube.com/watch?v=SXu89ikU4H0}} &
6.8 / 9.3 / 7.9 &
$\checkmark$ Много слов-паразитов и просторечий: fillers + non-literary excessive. &
$\checkmark$ VLM: U1,U3,U4; 3/5 кадров. Плотные записи на чёрной доске, визуальная неструктурированность. &
--- &
--- &
2 / 2 &
Студийный пример с большой чёрной доской: проверка, что высокое техническое качество может сочетаться с речевыми и визуальными замечаниями. \\

\midrule

\hypertarget{vid:v06}{V06} &
\textit{ПАРАМЕТР на ЕГЭ, который не решили даже ПРЕПОДЫ!}. Профиматика. {\tiny\url{https://www.youtube.com/watch?v=w-bOYpwIaa4}} &
4.0 / 8.0 / 5.3 &
$\checkmark$ Низкое качество звука: audio=4.0<5. &
$\checkmark$ VLM: F1,F3,F4,U2,U4; 4/5 кадров. Перегруженная доска-слайд, формулы и цветные элементы. Монтаж: jerky/min=0.72. &
--- &
$\checkmark$ Повышенный эмоциональный напор: high arousal, mean=0.772.&
2 / 4 &
Ярко инструментальный разбор параметра ЕГЭ; проверка tutorial-профиля, резкого монтажа, сенсорного флага и неформальной школьной подачи. \\

\midrule

\hypertarget{vid:v07}{V07} &
\textit{Введение в математический анализ 1. Множества и функции}. Лекторий ФПМИ. {\tiny\url{https://www.youtube.com/watch?v=SPMRNnpN-d8}} &
5.2 / 7.5 / 6.1 &
--- &
$\checkmark$ VLM: U3,U4; 3/5 кадров. Аудитория, следы на доске, фон/пересвет, провод в кадре. &
--- &
--- &
4 / 2 &
Длинная вводная лекция МФТИ с сильной переэкспозицией и академическим форматом; проверка визуальной устойчивости и профиля курса. \\

\midrule

\hypertarget{vid:v08}{V08} &
\textit{Введение в математический анализ 1. Действительные числа}. Лекторий ФПМИ. {\tiny\url{https://www.youtube.com/watch?v=uuesXwc1DJY}} &
4.0 / 7.1 / 5.1 &
$\checkmark$ Низкое качество звука: audio=4.0<5. &
$\checkmark$ VLM: U1,U3,U4; 2/5 кадров. Доска, провода, лишние предметы в аудитории. &
--- &
$\checkmark$ Повышенный эмоциональный напор: high arousal, mean=0.805. &
4 / 1 &
Тот же академический курс, но с очень плохим звуком при приемлемом видео; проверка общей технической агрегации и логического флага по аудио. \\

\midrule

\hypertarget{vid:v09}{V09} &
\textit{Ася Казанцева ``Зачем нужно спать и как делать это правильно?''}. Фонд ЗОЖ. {\tiny\url{https://www.youtube.com/watch?v=a0ftuZpigug}} &
8.0 / 5.0 / 6.2 &
--- &
$\checkmark$ VLM: F1,F2,F3; 3/5 кадров. Пёстрые слайды, коллажи, плотный текст и визуально шумный ряд. &
--- &
$\checkmark$ Низкая вариативность эмоциональной динамики: std=0.026.&
3 / 2 &
Научно-популярная лекция с дефектом произношения и средним визуальным качеством; проверка нарратива, sensor/VLM и границ детекции артикуляции. \\

\midrule

\hypertarget{vid:v10}{V10} &
\textit{Математический анализ. Превращение функций}. Институт мировой экономики и финансов. {\tiny\url{https://www.youtube.com/watch?v=SEQ3vqeO30k}} &
0.5 / 0.7 / 0.6 &
$\checkmark$ Очень низкое качество звука: audio=0.5<5; много просторечных/нестандартных речевых конструкций. &
$\checkmark$ VLM: U1,U4,F1; 3/5 кадров; video=0.7<5. Плохая картинка, доска, артефакты и элементы слайда. &
$\checkmark$ Отсутствуют смысловая динамика, образность и выраженный авторский голос.&
--- &
2 / 4 &
Крайний плохой случай по аудио/видео, но ярко инструментальный практический разбор; проверка нижней границы технического качества и instrumentality. \\

\midrule

\hypertarget{vid:v11}{V11} &
\textit{Лекция №2 по БЖД для студентов ФДО}. Кафедра ФК и ОБЖ. {\tiny\url{https://www.youtube.com/watch?v=keFFbjM2Zuc}} &
4.6 / 3.6 / 4.0 &
$\checkmark$ Низкое качество звука: audio=4.6<5.&
$\checkmark$ VLM: F1,F2,F3,F4; 5/5 кадров; video=3.6<5. Перегруженные слайды, вычурные фоны, плотный текст и цвета. &
$\checkmark$ Отсутствуют смысловая динамика, образность и выраженный авторский голос. &
$\checkmark$ Монотонная подача: mean=0.486, std=0.031.&
4 / 2 &
Крайний пример скучной лекции по слайдам с плохой презентацией; ожидаемый стресс-тест всех флагов, особенно sensor и intuitive. \\

\midrule

\hypertarget{vid:v12}{V12} &
\textit{Оземпик — лекарство от ожирения}. Александр Панчин. {\tiny\url{https://youtu.be/HOyEqwREGvw?si=4T_nngbvPjN2X_I3}} &
5.1 / 5.5 / 5.3 &
--- &
$\checkmark$ VLM: U1,U4; 3/5 кадров. Студийный фон: предметы, асимметрия. Монтаж: cuts/min=6.2. &
--- &
$\checkmark$ Низкая вариативность эмоциональной динамики: std=0.023. &
2 / 1 &
Пример студийного фона, который не должен автоматически считаться плохим; проверка VLM на ложные срабатывания и живого монтажа. \\

\midrule

\hypertarget{vid:v13}{V13} &
\textit{Владимир Жириновский — Лекция МГОУ (1999 г.)}. LDPR-TV. {\tiny\url{https://www.youtube.com/watch?v=FeWuTbDHE7g}} &
0.5 / 3.0 / 0.9 &
$\checkmark$ Очень низкое качество звука: audio=0.5<5. &
$\checkmark$ video=3.0<5. VLM не основной: проблема старой записи, а не перегруженности кадра. &
--- &
$\checkmark$ Высокий эмоциональный напор и активная жестикуляция: mean=0.838. &
1 / 1 &
Крайний пример эмоционального напора, старого видео и активной жестикуляции; проверка emotion-флага и технического штрафа старой записи. \\

\midrule

\hypertarget{vid:v14}{V14} &
\textit{Курс лекций ``Математический анализ''. Часть 12: Задачи на производные. Дифференциалы}. Российская экономическая школа. {\tiny\url{https://www.youtube.com/watch?v=7Pa53RISFak}} &
7.4 / 4.8 / 5.8 &
--- &
$\checkmark$ VLM: U1,U2,U4; 4/5 кадров; video=4.8<5. Фон, штора, провода и предметы у доски. &
--- &
--- &
4 / 3 &
Подготовительный курс с задачами: должен быть инструментальнее МФТИ, но сохранять академичность; проверка профиля «теория + практика». \\

\midrule

\hypertarget{vid:v15}{V15} &
\textit{СТРАДАНИЯ ИНОСТРАНЦА: почему русский язык такой сложный?}. Skyeng. {\tiny\url{https://www.youtube.com/watch?v=Ml_qgDhKC3w}} &
5.0 / 9.3 / 6.5 &
--- &
$\checkmark$ Монтаж: cuts/min=13.4, jerky/min=1.02. VLM не сработал. &
--- &
$\checkmark$ Активная жестикуляция.&
1 / 1 &
Пример акцента, студийного оформления и очень живого монтажа; важно, что фон не должен штрафоваться VLM, а sensor может сработать по склейкам. \\

\midrule

\hypertarget{vid:v16}{V16} &
\textit{Теория вероятности Дымов А.В. Лекция 06.09.24}. Матфак ВШЭ 22--26. {\tiny\url{https://youtu.be/chuYrk1oiaQ?si=1NwdNSUjpCqCzDcp}} &
8.3 / 9.6 / 8.9 &
$\checkmark$ Много слов-паразитов: fillers excessive.&
--- &
--- &
--- &
4 / 2 &
Крепкий академический курс по теории вероятностей для старших студентов; проверка высокого технического качества и стабильности академического профиля. \\

\bottomrule
\end{longtable}

\end{landscape}
\endgroup